% !TEX root = ../../M_1.tex
\subsection{Surds and indices}

Working with roots, powers of numbers, and algebraic expressions.

\subsubsection{Square, cube and higher roots}

\defn{Square root}{
  A \textbf{square root} of a number is a number that, when raised to the power of $2$, gives the original number. For example, $5$ and $-5$ are square roots of $25$ because $5^2 = 25$ and $(-5)^2 = 25$.
}

\defn{Cube root}{
  A \textbf{cube root} of a number is a number that, when raised to the power of $3$, gives the original number. For example, $3$ is a cube root of $27$ because $3^3 = 27$.
}

Fourth roots, fifth roots, and so on can be defined in a similar way. The $n$th root of a number $x$ is a number $y$ such that $y^n = x$.

\defn{Negative numbers}{
  Negative numbers do not have real square roots. For example, there is no real number $x$ such that $x^2 = -4$. However, negative numbers \textbf{do have real cube roots}. For example, $-3$ is a cube root of $-27$ because $(-3)^3 = -27$.
}

For all \textbf{even roots}, the root of a negative number is not a real number. For all \textbf{odd roots}, the root of a negative number is a negative number.

You can use the $\pm$ symbol to indicate that a number has two roots. For example, the square roots of $16$ are $\pm 4$, which means that both $4$ and $-4$ are square roots of $16$.

\prop{Using $\pm$ and $\sqrt{x}$}{
  Because $\sqrt{x}$ is used to denote a positive square root, we can write $\sqrt{16} = 4$. To indicate both roots, we can use $\pm$; for example, if $x^2 = 36$ then $x = \pm \sqrt{36} = \pm 6$.
}

The number to the upper left of the root symbol denotes the $n$th root. For example, $\sqrt[3]{64} = 4$ because $4^3 = 64$. The cube root of $-64$ is $-4$ because $(-4)^3 = -64$.

Most square roots are irrational numbers. For example, $\sqrt{3}$ is an irrational number because it cannot be expressed as a fraction of two integers. However, some square roots are rational numbers, such as $\sqrt{9} = 3$. \textbf{They are often just left in root form, such as $\sqrt{3}$, because they cannot be simplified.}

\defn{Surds}{
  An expression such as $3 + \sqrt{7}$ is called a \textbf{surd}. A surd is an expression that contains a root that cannot be simplified to remove the root. For example, $\sqrt{3}$ is a surd because it cannot be simplified to remove the square root. However, $\sqrt{9}$ is not a surd because it can be simplified to $3$.
}

Some common surds:

\[ \sqrt{2} \qquad \sqrt{3} \qquad \sqrt{5} \qquad \sqrt{7} \qquad \sqrt{11} \qquad \sqrt{13} \]

\subsubsection{Surd arithmetic}

You can manipulate surds using the normal rules of algebra and treat them the same way you would a variable, for example when collecting like terms.

Such as:

\[
  2\sqrt{5} - \sqrt{2} + 3\sqrt{5} + 4\sqrt{2} = (2 + 3)\sqrt{5} + (-1 + 4)\sqrt{2} = 5\sqrt{5} + 3\sqrt{2}
\]

\defn{Square root}{
  The \textbf{square root} $\sqrt{a}$ of a non-negative number $a$ is the non-negative number that gives $a$ when it is squared.
}

This definition is the whole reason the next rule works. Squaring and square-rooting undo each other, in the same way that multiplying by $2$ and dividing by $2$ undo each other.

\prop{Squaring a square root}{
  For any non-negative number $a$,
  \[ \sqrt{a} \times \sqrt{a} = \left(\sqrt{a}\right)^{2} = a. \]
  Multiplying a square root by itself removes the root.
}

A surd behaves like a variable, except at this step. For a variable, $x \times x = x^{2}$ and there is nothing more to do. For a surd, $\sqrt{6} \times \sqrt{6} = \left(\sqrt{6}\right)^{2}$ collapses to the whole number $6$.

\wex{a root multiplied by itself}{
  \begin{parts}
    \item Directly from the rule:
      \[ \sqrt{6} \times \sqrt{6} = 6, \qquad \left(\sqrt{11}\right)^{2} = 11. \]
      As a check, $\sqrt{6} \approx 2.449$ and $2.449 \times 2.449 \approx 6$.
    \item With a coefficient, multiply the numbers and the roots separately:
      \[ 5\sqrt{3} \times \sqrt{3} = 5 \times \left(\sqrt{3}\right)^{2} = 5 \times 3 = 15. \]
    \item When the whole term is squared, both the number and the root are squared:
      \[ \left(3\sqrt{5}\right)^{2} = 3^{2} \times \left(\sqrt{5}\right)^{2} = 9 \times 5 = 45. \]
    \item Inside a longer expression, the rule applies term by term:
      \[ \sqrt{2}\left(\sqrt{2} - 3\right) = \left(\sqrt{2}\right)^{2} - 3\sqrt{2} = 2 - 3\sqrt{2}. \]
    \item The rule is what removes a root from a denominator. Multiplying the top and bottom by $\sqrt{3}$ does not change the value:
      \[ \frac{1}{\sqrt{3}} = \frac{1 \times \sqrt{3}}{\sqrt{3} \times \sqrt{3}} = \frac{\sqrt{3}}{3}. \]
  \end{parts}
}

Some additional rules for manipulating surds:

\prop{Roots of products and quotients}{
  For non-negative numbers $a$ and $b$ (with $b \neq 0$ for the quotient),
  \[ \sqrt{ab} = \sqrt{a}\sqrt{b}, \qquad \sqrt{\frac{a}{b}} = \frac{\sqrt{a}}{\sqrt{b}}. \]
  The same rules hold for cube roots, fourth roots and so on.
}

The product rule follows from $\left(\sqrt{a}\right)^{2} = a$. Squaring $\sqrt{a}\sqrt{b}$ gives
\[ \left(\sqrt{a}\sqrt{b}\right)^{2} = \left(\sqrt{a}\right)^{2}\left(\sqrt{b}\right)^{2} = ab. \]
So $\sqrt{a}\sqrt{b}$ is a non-negative number whose square is $ab$, which is exactly what $\sqrt{ab}$ means.

\wex{combining and splitting roots}{
  \begin{parts}
    \item Combining two roots into one can give a whole number:
      \[ \sqrt{3}\sqrt{12} = \sqrt{3 \times 12} = \sqrt{36} = 6. \]
    \item The quotient rule works in the same way:
      \[ \frac{\sqrt{50}}{\sqrt{2}} = \sqrt{\frac{50}{2}} = \sqrt{25} = 5. \]
    \item Splitting a fraction under a root:
      \[ \sqrt{\frac{16}{25}} = \frac{\sqrt{16}}{\sqrt{25}} = \frac{4}{5}. \]
    \item Cube roots follow the same pattern:
      \[ \sqrt[3]{3}\,\sqrt[3]{9} = \sqrt[3]{3 \times 9} = \sqrt[3]{27} = 3. \]
    \item Used in reverse, the product rule simplifies a surd. Split the number into a product where one factor is a perfect square:
      \[ \sqrt{45} = \sqrt{9 \times 5} = \sqrt{9}\sqrt{5} = 3\sqrt{5}. \]
  \end{parts}
}

The rule does not work for addition: $\sqrt{a + b} \neq \sqrt{a} + \sqrt{b}$. For example:
\[ \sqrt{36 + 64} = \sqrt{100} = 10, \quad \text{but} \quad \sqrt{36} + \sqrt{64} = 6 + 8 = 14. \]

Another way to simplify surds applies if the number inside the square root sign has a factor that is a perfect square. For example, $\sqrt{72}$ can be simplified because $72 = 36 \times 2$ and $36$ is a perfect square. This is due to the fact that:

\[
\sqrt{ab} = \sqrt{a}\sqrt{b}
\]

\wex{pulling out square factors}{
  Consider the following surds: $\sqrt{80}$ and $\sqrt{15}$.
  \begin{parts}
    \item $80$ has a factor of $16$, which is a perfect square. So we can write:
      \[ \sqrt{80} = \sqrt{16 \times 5} = \sqrt{16}\sqrt{5} = 4\sqrt{5}. \]
    \item $15$ has no factors that are perfect squares (other than $1$), so it cannot be simplified. So we write:
      \[ \sqrt{15} = \sqrt{15}. \] As it can't be simplified, it is left in root form.
  \end{parts}
}

When simplifying a surd, look for the largest perfect square that is a factor of the number under the root. The first 15 perfect squares are
\[ 1,\ 4,\ 9,\ 16,\ 25,\ 36,\ 49,\ 64,\ 81,\ 100,\ 121,\ 144,\ 169,\ 196,\ 225. \]

\defn{Rationalising the denominator}{
  Rewriting a fraction so that its denominator no longer contains an irrational root is called \textbf{rationalising} the denominator. It is done by multiplying the top and bottom of the fraction by a suitable surd, which does not change the value of the fraction.
}

When the denominator is a single root, multiply the top and bottom by that root. The rule $\left(\sqrt{a}\right)^{2} = a$ then removes the root from the denominator:
\[ \frac{1}{\sqrt{5}} = \frac{1}{\sqrt{5}} \times \frac{\sqrt{5}}{\sqrt{5}} = \frac{\sqrt{5}}{5}. \]

\defn{Conjugate}{
  A \textbf{conjugate} of a two-term expression is the expression obtained by changing the sign of one of its terms, usually the second. For example, a conjugate of $2\sqrt{3} + \sqrt{5}$ is $2\sqrt{3} - \sqrt{5}$, and a conjugate of $5 - 3\sqrt{2}$ is $5 + 3\sqrt{2}$.
}

\prop{Strategy for rationalising the denominator}{
  \begin{itemize}[topsep=2pt, itemsep=2pt]
    \item If the denominator is a single term containing $\sqrt{a}$, multiply the top and bottom by $\sqrt{a}$.
    \item If the denominator has two terms involving square roots, multiply the top and bottom by a conjugate of the denominator.
  \end{itemize}
}

The conjugate works because the new denominator has the form $(A + B)(A - B)$, which is the difference of two squares $A^{2} - B^{2}$. Squaring each term removes its square root.

\wex{clearing roots from a denominator}{
  \begin{parts}
    \item Only the $\sqrt{5}$ needs removing, so multiply by $\sqrt{5}$:
      \[ \frac{2}{7\sqrt{5}} = \frac{2}{7\sqrt{5}} \times \frac{\sqrt{5}}{\sqrt{5}} = \frac{2\sqrt{5}}{7 \times 5} = \frac{2\sqrt{5}}{35}. \]
    \item A conjugate of $5 - 3\sqrt{2}$ is $5 + 3\sqrt{2}$:
      \begin{align*}
        \frac{\sqrt{3}}{5 - 3\sqrt{2}}
          &= \frac{\sqrt{3}}{5 - 3\sqrt{2}} \times \frac{5 + 3\sqrt{2}}{5 + 3\sqrt{2}}
           = \frac{5\sqrt{3} + 3\sqrt{2}\sqrt{3}}{5^{2} - \left(3\sqrt{2}\right)^{2}} \\[4pt]
          &= \frac{5\sqrt{3} + 3\sqrt{6}}{25 - 9 \times 2}
           = \frac{5\sqrt{3} + 3\sqrt{6}}{7}
      \end{align*}
  \end{parts}
}

A final answer that is a surd should be simplified as far as possible. Rationalising the denominator is not always required, so leaving $1/\sqrt{5}$ as it is is acceptable. Many calculators rationalise automatically and display $1/\sqrt{5}$ as $\sqrt{5}/5$.

\subsubsection{Powers and the index laws}

\defn{Base and index}{
  In $3^{4} = 3 \times 3 \times 3 \times 3 = 81$, the number $3$ is the \textbf{base} and $4$ is the \textbf{power}, \textbf{index} or \textbf{exponent}. The plural of index is \textbf{indices}. The result, $81$, is also called a power of $3$.
}

\prop{Index laws for a single base}{
  \[ a^{m}a^{n} = a^{m + n}, \qquad \frac{a^{m}}{a^{n}} = a^{m - n}, \qquad \left(a^{m}\right)^{n} = a^{mn}. \]
  The first law extends to more than two powers, for example $a^{m}a^{n}a^{r} = a^{m + n + r}$.
}

When considering \textbf{negative indices}, such as $3^{-3}$, we can see that this is the same as the \textbf{reciprocal} of $3^{3}$, which is $\frac{1}{3^{3}} = \frac{1}{27}$.

\prop{Zero and negative index laws for a single base}{
  For any non-zero number $a$:
  \begin{itemize}[topsep=2pt, itemsep=2pt]
    \item Any power with index $0$ equals $1$:
      \[ a^{0} = 1. \]
    \item A negative index means taking the reciprocal of the matching positive power:
      \[ a^{-n} = \frac{1}{a^{n}}. \]
    \item The reverse also holds, so a negative power in a denominator becomes a positive power in the numerator:
      \[ \frac{1}{a^{-n}} = a^{n}, \]
      since $\dfrac{1}{a^{-n}} = 1 \div \dfrac{1}{a^{n}} = a^{n}$.
  \end{itemize}
  With $n = 1$, this gives $a^{-1} = \dfrac{1}{a}$, so an index of $-1$ simply gives the reciprocal.
}

The index laws above involve just one base number, but there are two more index laws that involve two base numbers. 


\prop{Index laws for two base numbers}{
  The powers of a product and a quotient can be expressed as follows:
      \[ (ab)^{n} = a^{n}b^{n}, \quad \left(\frac{a}{b}\right)^{n} = \frac{a^{n}}{b^{n}}. \]
}

Working this through an example:

\wex{negative indices as a product of fractions}{
  Write each factor as its own fraction. A factor with a negative index flips to the other side of its fraction line and its index becomes positive. Then multiply the fractions back together.
  \begin{parts}
    \item Only the denominator has a negative index:
      \begin{align*}
        \frac{x^{3}}{x^{-2}}
          &= \frac{x^{3}}{1} \times \frac{1}{x^{-2}}
           = \frac{x^{3}}{1} \times \frac{x^{2}}{1} \\[4pt]
          &= x^{3}x^{2}
           = x^{5}
      \end{align*}
    \item Raise every factor in the term to the power $3$ first, then flip the factor with the negative index:
      \begin{align*}
        \left(3x^{-2}p\right)^{3}
          &= 3^{3} \times \left(x^{-2}\right)^{3} \times p^{3}
           = \frac{27}{1} \times \frac{x^{-6}}{1} \times \frac{p^{3}}{1} \\[4pt]
          &= \frac{27}{1} \times \frac{1}{x^{6}} \times \frac{p^{3}}{1}
           = \frac{27p^{3}}{x^{6}}
      \end{align*}
  \end{parts}

  \textbf{Index laws used.}
  In (a), $\dfrac{1}{a^{-n}} = a^{n}$ flips $x^{-2}$ into the numerator, then $a^{m}a^{n} = a^{m + n}$ combines $x^{3}x^{2}$ into $x^{5}$.
  In (b), $(ab)^{n} = a^{n}b^{n}$ applies the power to every factor, $\left(a^{m}\right)^{n} = a^{mn}$ turns $\left(x^{-2}\right)^{3}$ into $x^{-6}$, and $a^{-n} = \dfrac{1}{a^{n}}$ moves $x^{-6}$ into the denominator as $x^{6}$.
}

\subsubsection{Fractional powers}